% =====================================================================
%  A3 -- Economics: two-product profit maximisation
% =====================================================================
\subsection{A3: Economics, two-product profit maximisation}
\setprob{A3}

\begin{frame}{\probtitle{A3}{Maximising a two-product profit by gradient ascent}}
  \begin{probbox}[Problem A3 \hfill Application: Economics \quad $\star\star$]
    A firm sells two substitute products in quantities $x$ and $y$ (hundreds of units). Its profit,
    in thousand Tk, is
    \[
      P(x,y) = 100x + 80y - 2x^2 - 2y^2 - 2xy .
    \]
    Each period the firm adjusts production by \emph{gradient ascent}, $\x_{k+1} = \x_k + \alpha\nabla P(\x_k)$.
    \begin{enumerate}[(a)]
      \item Find the optimum analytically and show that it is a global maximum.
      \item Find the stable range of $\alpha$ and the best fixed $\alpha$.
      \item Run three iterations from $(0,0)$ with the best $\alpha$, and tabulate $P$ and the error.
      \item Explain what $\alpha = 0.4$ does, in business terms.
    \end{enumerate}
  \end{probbox}
  \footnotesize\textcolor{mGrey}{The point: ascent is just descent on $-P$, and a hand calculation matches the theory exactly.}
\end{frame}

\begin{frame}{\soltitle{A3}{1}{3}}
  \textbf{(a)}
  $\nabla P = \begin{pmatrix}100 - 4x - 2y\\ 80 - 2x - 4y\end{pmatrix} = \mathbf 0
  \iff \begin{cases}4x + 2y = 100\\ 2x + 4y = 80\end{cases}
  \iff \ans{(x^\star,y^\star) = (20,\ 10)}$
  \[
    \nabla^2P = \begin{pmatrix}-4&-2\\-2&-4\end{pmatrix}, \quad \text{eigenvalues } -2,\,-6
    \;\Rightarrow\; \text{negative definite} \;\Rightarrow\; \text{global maximum}
  \]
  \[
    P^\star = 2000 + 800 - 800 - 200 - 400 = \ans{1400} \text{ thousand Tk}.
  \]
  \textbf{(b)} Ascent on $P$ is GD on $-P$, with $\bA = \begin{pmatrix}4&2\\2&4\end{pmatrix}$:
  $\lambda = 2$ along $(1,-1)$ and $\lambda = 6$ along $(1,1)$.
  \[
    \ans{0 < \alpha < \tfrac26 = \tfrac13},\qquad
    \alpha^\star = \tfrac{2}{2+6} = \ans{\tfrac14},\qquad
    \rho = \tfrac{6-2}{6+2} = \ans{\tfrac12}.
  \]
\end{frame}

\begin{frame}{\soltitle{A3}{2}{3}}
  \begin{columns}[T]
    \begin{column}{0.56\textwidth}
      \textbf{(c)} $\alpha = \tfrac14$, starting from $(0,0)$:

      \vspace{2pt}
      \footnotesize
      \begin{tabular}{@{}c c c c c@{}}
        \toprule
        $k$ & $\x_k$ & $\nabla P(\x_k)$ & $P$ & $\norm{\e_k}$\\
        \midrule
        0 & $(0,0)$ & $(100,80)$ & $0$ & $22.361$\\
        1 & $(25,20)$ & $(-40,-50)$ & $1050$ & $11.180$\\
        2 & $(15,7.5)$ & $(25,20)$ & $1312.5$ & $5.590$\\
        3 & $(21.25,12.5)$ & $(-10,-12.5)$ & $1378.125$ & $2.795$\\
        \bottomrule
      \end{tabular}

      \vspace{4pt}
      \small
      \begin{itemize}
        \item The error \alert{halves exactly} every step, because both factors, $1 - \tfrac14\cdot2 = \tfrac12$ and
              $1 - \tfrac14\cdot6 = -\tfrac12$, have size $\tfrac12$.
        \item The profit gap shrinks by $\rho^2 = \tfrac14$ per step: $1400 \to 350 \to 87.5 \to 21.875$.
      \end{itemize}
    \end{column}
    \begin{column}{0.42\textwidth}
      \centering
      \includegraphics[width=0.95\linewidth]{fig_profit}
    \end{column}
  \end{columns}
\end{frame}

\begin{frame}{\soltitle{A3}{3}{3}}
  \textbf{The overshoot pattern.} The $(1,1)$ part of the error changes sign every step (its factor is $-\tfrac12$).
  So total production first \emph{overshoots} ($25+20 = 45 > 30$), then \emph{undershoots}, then overshoots again,
  by half as much each time.

  \smallskip
  \textbf{(d) $\alpha = 0.4 > \tfrac13$.} Now the factor along $(1,1)$ is $1 - 0.4\cdot6 = \alert{-1.4}$, and
  $\x_6 \approx (-92.9,\ -102.9)$. The swings grow instead of dying out.
  \smallskip
  In business terms: a management that \emph{over-reacts} to the marginal-profit signal pushes total output
  further from the optimum every period, until it is planning negative (meaningless) quantities.
  \begin{keybox}[Lesson]
    How fast you adjust has to respect the curvature of the profit function: $\alpha < 2/\lmax$.
    Because the goods are substitutes (the $-2xy$ term), $\lmax$ rises from $4$ to $6$ and the safe step gets smaller.
  \end{keybox}
\end{frame}
